\documentclass[11pt,a4paper]{article}
\usepackage[T1]{fontenc}
\usepackage{lmodern,amsmath,amssymb,amsthm,mathtools,microtype}
\usepackage[margin=27mm]{geometry}
\usepackage[colorlinks=true,linkcolor=blue,citecolor=blue,urlcolor=blue]{hyperref}
\setlength{\emergencystretch}{2em}
\newtheorem{theorem}{Theorem}[section]
\newtheorem{proposition}[theorem]{Proposition}
\newtheorem{corollary}[theorem]{Corollary}
\newtheorem{remark}[theorem]{Remark}
\newcommand{\Q}{\mathbb Q}
\newcommand{\Pp}{\mathbb P}
\newcommand{\Jac}{\operatorname{Jac}}
\newcommand{\Gal}{\operatorname{Gal}}
\newcommand{\divisor}{\operatorname{div}}
\title{Fixed-class synchronization of normalized pole parameters}
\author{}
\date{16 November 2026}
\begin{document}
\maketitle
\begin{abstract}
The simultaneous preservation of prescribed quadratic classes in the
functions $(t^2-D)/(t^2+3D)$ defines an explicit multiquadratic curve.
For distinct positive parameters its genus is computed, and in the
two-parameter case its Jacobian is decomposed into three elliptic factors
and a genus-two Jacobian. The common-parameter product of the associated
quartic families has a smooth proper model of Kodaira dimension one.
\end{abstract}

\section{Fixed quadratic propagation classes}

For $D\in\Q_{>0}$ set
\[
 e_D(t)=\frac{t^2-D}{t^2+3D},\qquad
 F_D(t)=(t^2-D)(t^2+3D).
\]
All curves below mean smooth projective normalizations of the displayed
affine equations.

The quadratic class constraint is the one produced by a specific
propagation rule. If $(t_0,s_0,q_0)$ satisfies
\[
 q_0^2=D(s_0^2-Ce_D(t_0))(s_0^2+3Ce_D(t_0))
\]
and $e_D(t)=e_D(t_0)z^2$, then $(s,q)=(zs_0,z^2q_0)$ satisfies the
same pole equation at $t$. The theorems below concern simultaneous
use of this rule with fixed classes. They do not constrain realizations
that vary those classes.

\begin{theorem}\label{thm:genus}
Let $D_1,\ldots,D_r$ be distinct positive rational numbers and let
$\alpha_1,\ldots,\alpha_r\in\Q^\times$. The curve
\[
 B:\qquad \alpha_i z_i^2=e_{D_i}(t)\quad(1\leq i\leq r)
\]
is geometrically integral. Its map to $\Pp^1_t$ has degree $2^r$, and
\[
                 g(B)=1+2^r(r-1).
\]
\end{theorem}
\begin{proof}
Over an algebraic closure each function $e_{D_i}$ has odd valuation at
exactly four points, namely $t^2=D_i$ and $t^2=-3D_i$. These sets are
pairwise disjoint: equal signs require equal $D_i$, and opposite signs
cannot occur because the $D_i$ are positive. Any nonempty product of the
$e_{D_i}$ has an odd valuation at a branch point of one selected factor.
The corresponding square classes are therefore independent over
$\overline{\Q}(t)$. This proves geometric integrality and degree $2^r$.

At each of the $4r$ branch points the inertia group has order two; the
other extensions are unramified and split over the geometric completed
local field. The ramification contribution at such a point is
$2^{r-1}$. Riemann--Hurwitz gives
\[
 2g(B)-2=-2\cdot2^r+4r\cdot2^{r-1}=2^{r+1}(r-1).
\]
\end{proof}

\begin{corollary}
For $r\geq2$, each fixed tuple of classes $\alpha_i$ admits only finitely
many common rational parameter values. In particular, any finite union
of fixed tuples has this property.
\end{corollary}
\begin{proof}
The genus in Theorem~\ref{thm:genus} is greater than one. Apply the
finiteness theorem for rational points on curves of genus greater than
one~\cite{Faltings}.
\end{proof}

\begin{remark}
For two factors the genus is five, and for three it is seventeen.
Finiteness for a fixed tuple of classes does not exclude a common
rational parameter, and it does not control a union over infinitely many
class tuples.
\end{remark}

\begin{proposition}\label{prop:mobius}
Let $D_1,D_2>0$ and suppose $D_1/D_2$ is not a rational square. If
$\phi\in\operatorname{PGL}_2(\Q)$, the normalization of
\[
 \alpha_1z_1^2=e_{D_1}(t),\qquad
 \alpha_2z_2^2=e_{D_2}(\phi(t))
\]
has genus five.
\end{proposition}
\begin{proof}
A branch point of the first map has residue field
$\Q(\sqrt{D_1})$ or $\Q(\sqrt{-3D_1})$. Rational fractional linear
transformations preserve residue fields, because they and their inverses
are defined over $\Q$. The corresponding fields for the second map are
$\Q(\sqrt{D_2})$ and $\Q(\sqrt{-3D_2})$.
A real field cannot equal an imaginary quadratic field. Within the real
pair or the imaginary pair, equality is equivalent to $D_1/D_2$ being a
square, including the possibility that a real residue field is $\Q$.
Thus the eight geometric branch points are distinct. The proof of
Theorem~\ref{thm:genus} applies with $r=2$.
\end{proof}

\section{The genus-five quotients}

\begin{proposition}\label{prop:decomposition}
For $r=2$ write $w_i=z_i(t^2+3D_i)$, so that
\[
                    \alpha_iw_i^2=F_{D_i}(t).
\]
Let $B_i$ denote the two individual genus-one curves. Define
\begin{align*}
 E_{12}:\quad \alpha_1\alpha_2 Z^2
   &=(u-D_1)(u+3D_1)(u-D_2)(u+3D_2),\\
 H_{12}:\quad \alpha_1\alpha_2 V^2
   &=u(u-D_1)(u+3D_1)(u-D_2)(u+3D_2).
\end{align*}
Then $E_{12}$ has genus one, $H_{12}$ has genus two, and over $\Q$ there
is an isogeny
\[
 \Jac(B)\ \sim\ \Jac(B_1)\times\Jac(B_2)
                  \times\Jac(E_{12})\times\Jac(H_{12}).
\]
\end{proposition}
\begin{proof}
The five finite branch values $0,D_1,-3D_1,D_2,-3D_2$ are distinct.
This gives the asserted genera. The three nontrivial character quotients
for the sign changes of $(w_1,w_2)$ are $B_1$, $B_2$, and
\[
 G:\quad \alpha_1\alpha_2 W^2=F_{D_1}(t)F_{D_2}(t),
 \qquad W=w_1w_2.
\]
Their genera are $1,1,3$, respectively. Pullback on Jacobians gives
\(\Jac(B_1)\times\Jac(B_2)\times\Jac(G)\to\Jac(B)\).
The pulled-back spaces of regular differentials are the three distinct
character spaces and have total dimension five. Hence this map is an
isogeny.

The genus-three curve $G$ maps to $E_{12}$ by
$(u,Z)=(t^2,W)$ and to $H_{12}$ by $(u,V)=(t^2,tW)$.
The respective involutions are $(t,W)\mapsto(-t,W)$ and
$(t,W)\mapsto(-t,-W)$. Their regular differential spaces have dimensions
one and two and exhaust that of $G$. The same argument gives
$\Jac(G)\sim\Jac(E_{12})\times\Jac(H_{12})$.
\end{proof}

The genus-one quotient by $t\mapsto-t$ can also be written
\[
 D_1\frac{1+3\alpha_1z_1^2}{1-\alpha_1z_1^2}
 =D_2\frac{1+3\alpha_2z_2^2}{1-\alpha_2z_2^2}=u.
\]
A point on this quotient lifts only if $u$ is a rational square.
The genus-two factor in Proposition~\ref{prop:decomposition} is an
additional explicit part of the compatibility, rather than a condition
that can be omitted.

\section{The common-parameter variety}

\begin{theorem}\label{thm:canonical}
Let $D_1,\ldots,D_r$ be distinct positive rational numbers, and let
$C_1,\ldots,C_r\in\Q^\times$. The variety
\[
 q_i^2=D_i\bigl(s_i^2-C_ie_{D_i}(t)\bigr)
               \bigl(s_i^2+3C_ie_{D_i}(t)\bigr)
                  \quad(1\leq i\leq r)
\]
has a smooth proper model $f:X\to\Pp^1_t$ with
\[
                   \omega_X\simeq f^*\mathcal O_{\Pp^1}(2r-2).
\]
For $r\geq2$ its Kodaira dimension is one, and for $m\geq1$
\[
                 h^0(X,\omega_X^{\otimes m})=(2r-2)m+1.
\]
\end{theorem}
\begin{proof}
Let $B$ be the curve of Theorem~\ref{thm:genus} with all $\alpha_i=1$,
and let $A_i$ be the smooth genus-one curve
\[
                  v_i^2=D_i(u_i^2-C_i)(u_i^2+3C_i).
\]
The product $Y=B\times\prod_i A_i$ has an action of
$G=(\mathbb Z/2\mathbb Z)^r$. Generator $i$ changes the signs of
$z_i$ and $u_i$, leaving the other coordinates fixed. The invariant
functions
\[
                  s_i=z_i u_i,\qquad q_i=z_i^2v_i
\]
identify the function field of $Y/G$ with the function field in the
statement. Indeed, after adjoining all $z_i$, these formulas recover
$u_i,v_i$, and the resulting extension has the displayed faithful group
of order $2^r$.

Away from the branch points of $B\to\Pp^1$ the group acts freely on
$B$. At a branch point its stabilizer is generated by exactly one of
the generators. On $A_i$, the involution $u_i\mapsto-u_i$ has four
fixed points and has derivative $-1$ at each. Consequently every
nontrivial stabilizer on $Y$ has order two and acts on tangent
coordinates as
\[
                     (-1,-1,1,\ldots,1).
\]
The quotient singularities are products of a surface ordinary double
point with a smooth variety. Their singular loci are disjoint. Resolving
the transverse ordinary double points gives a smooth proper model $X$,
and the resolution is crepant. Indeed, in local transverse coordinates
$(\xi,\eta)$ the invariant coordinates
$a=\xi^2$, $b=\eta^2$, $c=\xi\eta$ satisfy $ab=c^2$. The canonical
form $da\wedge dc/a$ pulls back to $2\,d\xi\wedge d\eta$.
On the resolution chart $c=av$ it becomes $da\wedge dv$, and the
other chart is identical with $a,b$ interchanged. Thus the discrepancy
is zero.

Put $w_i=z_i(t^2+3D_i)$. On $Y$ consider the rational top form
\[
 \Omega=\frac{dt}{w_1\cdots w_r}
                    \wedge\frac{du_1}{v_1}\wedge\cdots
                    \wedge\frac{du_r}{v_r}.
\]
Each $du_i/v_i$ is regular and nowhere zero. At a finite branch point,
$dt$ and the unique vanishing $w_i$ each have order one; all remaining
$w_j$ are units. Thus $\Omega$ has no finite zeros or poles. At each
point above infinity, the map $B\to\Pp^1$ is unramified,
$\operatorname{ord}(dt)=-2$ and $\operatorname{ord}(w_i)=-2$.
It follows that
\[
               \divisor_Y(\Omega)=(2r-2)(Y\to\Pp^1)^*(\infty).
\]
Every generator changes the signs of one $w_i$ and one differential
$du_i/v_i$, so $\Omega$ is invariant. The quotient is unramified in
codimension one, and its crepant resolution does not change the
canonical divisor. The form therefore descends and gives
$\omega_X\simeq f^*\mathcal O(2r-2)$.

The geometric generic fiber is a product of smooth genus-one curves,
so it is geometrically connected. Stein factorization yields
$f_*\mathcal O_X=\mathcal O_{\Pp^1}$. The projection formula now gives
the claimed plurigenera and Kodaira dimension.
\end{proof}

\begin{corollary}
For $r\geq2$, the common-parameter variety is not birational to a
quotient of an abelian variety by a finite group.
\end{corollary}
\begin{proof}
A generically finite dominant rational map from an abelian variety to
$X$ would imply $\kappa(X)\leq0$, by pulling back pluricanonical forms
after resolution. This contradicts Theorem~\ref{thm:canonical}.
\end{proof}


\section{Exact algebra checks}

The branch calculations can be verified by
\[
 \operatorname{disc}_t F_D=-12288D^6,
\]
\[
 \operatorname{Res}_t(F_{D_1},F_{D_2})
 =9(D_1-D_2)^4(D_1+3D_2)^2(3D_1+D_2)^2.
\]
The accompanying script
\begin{center}\small
\nolinkurl{negative_results/computation/}\\
\nolinkurl{fixed_class_synchronization_v43/verify.py}
\end{center}
checks these identities, the scaling identities, the local quotient
relation, and the numerical Riemann--Hurwitz and character dimensions.
The finiteness and canonical-bundle conclusions follow from the proofs
above, rather than from a finite rational-point search.

\begin{thebibliography}{9}
\bibitem{Faltings}
G. Faltings, \emph{Endlichkeitss\"atze f\"ur abelsche Variet\"aten
\"uber Zahlk\"orpern}, Invent. Math. \textbf{73} (1983), 349--366,
doi:10.1007/BF01388432; correction, \textbf{75} (1984), 381.
\end{thebibliography}
\end{document}
